% ============================================================================
%  headermeyer-prob-stats.tex — shared preamble for the Meyer solutions manual
%
%  Included by master.tex and by each of the ten per-chapter documents via:
%      \IfFileExists{headermeyer-prob-stats.tex}{% ============================================================================
%  headermeyer-prob-stats.tex — shared preamble for the Meyer solutions manual
%
%  Included by master.tex and by each of the ten per-chapter documents via:
%      \IfFileExists{headermeyer-prob-stats.tex}{% ============================================================================
%  headermeyer-prob-stats.tex — shared preamble for the Meyer solutions manual
%
%  Included by master.tex and by each of the ten per-chapter documents via:
%      \IfFileExists{headermeyer-prob-stats.tex}{% ============================================================================
%  headermeyer-prob-stats.tex — shared preamble for the Meyer solutions manual
%
%  Included by master.tex and by each of the ten per-chapter documents via:
%      \IfFileExists{headermeyer-prob-stats.tex}{\input{headermeyer-prob-stats}}{\input{../headermeyer-prob-stats}}
%  (the fallback lets a chapter compile whether the working directory is the
%   repository root or the chapter's own folder).
%
%  This file is the *union* of the ten original per-chapter preambles, so every
%  chapter keeps the packages, macros and theorem environments it relied on.
%  Two divergences in the originals are reconciled here:
%    * inputenc: chapter 1 declared utf8, the others latin1. Every source file
%      is in fact UTF-8, and the only non-ASCII byte in any of them is a "§"
%      inside a comment, so utf8 is used throughout and renders identically.
%    * theorem environments: chapters 7 and 8 used the starred names defn*,
%      cor*, prop*, thm*; the other chapters used defn, cor, prop, thm. Both
%      families are defined below, with their original numbering preserved.
% ============================================================================

\PassOptionsToPackage{dvipsnames}{xcolor}

\usepackage{geometry}
\geometry{letterpaper}
\geometry{
 a4paper,
 total={170mm,257mm},
 left=15mm,
 right = 15mm,
 top=15mm,
 bottom = 20mm
 }

\usepackage{graphicx}
% Figures live alongside the chapter that uses them. Listing "./" first plus
% every chapter folder lets a document find its figures whether it is compiled
% from inside its own folder or from the repository root (which is how the
% master builds), so the chapters no longer depend on the working directory.
\graphicspath{%
  {./}%
  {Chapter 1 - Introduction to Probability/}%
  {Chapter 2 - Finite Sample Spaces/}%
  {Chapter 3 - Conditional Probability and Independence/}%
  {Chapter 4 - One-Dimensional Random Variables/}%
  {Chapter 5 - Functions of Random Variables/}%
  {Chapter 6 - Two- and Higher-Dimensional Random Variables/}%
  {Chapter 7 - Further Characteristics of Random Variables/}%
  {Chapter 8 - The Poisson and Other Discrete Random Variables/}%
  {Chapter 9 - Some Important Continuous Random Variables/}%
  {Chapter 10 - The Moment-Generating Function/}%
}
\usepackage{setspace}
\setstretch{0.5}

\usepackage{lmodern}
\usepackage{amsmath,amssymb,amsthm,enumitem,mathtools,xpatch}
\usepackage{bm}
\usepackage[most]{tcolorbox}
\usepackage[dvipsnames]{xcolor}
\usepackage{float}
\usepackage{mathrsfs}
\usepackage{wrapfig, lipsum, amsthm, thmtools}
\usepackage{framed}
\usepackage{gensymb}
\usepackage{url}

\newcommand*{\simsym}{\mathord\sim}

\newcommand*{\Perm}[2]{{}^{#1}\!P_{#2}}%
\newcommand*{\Comb}[2]{{}^{#1}C_{#2}}%

\usepackage[framemethod=tikz]{mdframed}

% ---------------------------------------------------------------------------
% Theorem environments
% ---------------------------------------------------------------------------
\theoremstyle{definition}
\newtheorem*{exmp*}{Example}

% Family used by chapters 1-6, 9, 10
\newtheorem*{defn}{Definition}
\surroundwithmdframed[backgroundcolor=white]{defn}

\newtheorem{cor}{Corollary}
\surroundwithmdframed[backgroundcolor=white]{cor}

\newtheorem{prop}{Proposition}
\surroundwithmdframed[backgroundcolor=white]{prop}

\newtheorem*{thm}{Theorem}
\surroundwithmdframed[backgroundcolor=white]{thm}

% Family used by chapters 7 and 8
\newtheorem*{defn*}{Definition}
\surroundwithmdframed[backgroundcolor=white]{defn*}

\newtheorem{cor*}{Corollary}
\surroundwithmdframed[backgroundcolor=white]{cor*}

\newtheorem*{prop*}{Proposition}
\surroundwithmdframed[backgroundcolor=white]{prop*}

\newtheorem*{thm*}{Theorem}
\surroundwithmdframed[backgroundcolor=white]{thm*}

% ---------------------------------------------------------------------------
% Encoding
% ---------------------------------------------------------------------------
\usepackage[utf8]{inputenc}

% ---------------------------------------------------------------------------
% tikz (probability trees) and pgfplots
% ---------------------------------------------------------------------------
\usepackage{tikz}
\usetikzlibrary{trees,calc,angles,positioning,intersections}

\usepackage{pgfplots}
\pgfplotsset{width=10cm,compat=1.9}
\usepgfplotslibrary{fillbetween}

% quotations dirty talk
\usepackage{dirtytalk}

% floor ceiling
\DeclarePairedDelimiter\ceil{\lceil}{\rceil}
\DeclarePairedDelimiter\floor{\lfloor}{\rfloor}
\DeclarePairedDelimiter\abs{\lvert}{\rvert}
% plain \abs{x} autosizes like \left| x \right|; \abs*{x} keeps fixed-size bars
\DeclarePairedDelimiter\absfixed{\lvert}{\rvert}%
\RenewDocumentCommand{\abs}{s}{\IfBooleanTF{#1}{\absfixed}{\absfixed*}}%

% diag table
\usepackage{diagbox}

% math circle
\makeatletter
\newcommand\mathcircled[1]{%
  \mathpalette\@mathcircled{#1}%
}
\newcommand\@mathcircled[2]{%
  \tikz[baseline=(math.base)] \node[draw,circle,inner sep=1pt] (math) {$\m@th#1#2$};%
}
\makeatother

% ---------------------------------------------------------------------------
% Title block
% ---------------------------------------------------------------------------
% Follows the convention used in the deqla and apde manuals: the document
% defines \nbook (and \nbookshort) before inputting this header, and the title
% block is assembled here from those macros -- \title is left empty and the
% book, its author and the solutions author are stacked in \author, with a
% closing disclaimer appended to \maketitle.
%
% Guarded on \nbook so only the merged document picks this up; the ten chapter
% documents declare their own \title/\author after inputting this file and are
% left exactly as they were. Defining \nisofficial suppresses the disclaimer.
\ifx \nbook \undefined
\else
  \ifx \nauthor \undefined
    \def\nauthor{David A. Lee}
  \fi
  \author{\nbook \\ \small Paul L. Meyer \\ \small Solutions by \nauthor}
  \date{}
  \title{}
  \ifx \nisofficial \undefined
    \makeatletter
    \let\@real@maketitle\maketitle
    \renewcommand{\maketitle}{\@real@maketitle\begin{center}\begin{minipage}[c]{0.9\textwidth}\centering\footnotesize All errors, typographical and substantive, and other offenses, are entirely my own.\end{minipage}\end{center}}
    \makeatother
  \fi
\fi
}{\input{../headermeyer-prob-stats}}
%  (the fallback lets a chapter compile whether the working directory is the
%   repository root or the chapter's own folder).
%
%  This file is the *union* of the ten original per-chapter preambles, so every
%  chapter keeps the packages, macros and theorem environments it relied on.
%  Two divergences in the originals are reconciled here:
%    * inputenc: chapter 1 declared utf8, the others latin1. Every source file
%      is in fact UTF-8, and the only non-ASCII byte in any of them is a "§"
%      inside a comment, so utf8 is used throughout and renders identically.
%    * theorem environments: chapters 7 and 8 used the starred names defn*,
%      cor*, prop*, thm*; the other chapters used defn, cor, prop, thm. Both
%      families are defined below, with their original numbering preserved.
% ============================================================================

\PassOptionsToPackage{dvipsnames}{xcolor}

\usepackage{geometry}
\geometry{letterpaper}
\geometry{
 a4paper,
 total={170mm,257mm},
 left=15mm,
 right = 15mm,
 top=15mm,
 bottom = 20mm
 }

\usepackage{graphicx}
% Figures live alongside the chapter that uses them. Listing "./" first plus
% every chapter folder lets a document find its figures whether it is compiled
% from inside its own folder or from the repository root (which is how the
% master builds), so the chapters no longer depend on the working directory.
\graphicspath{%
  {./}%
  {Chapter 1 - Introduction to Probability/}%
  {Chapter 2 - Finite Sample Spaces/}%
  {Chapter 3 - Conditional Probability and Independence/}%
  {Chapter 4 - One-Dimensional Random Variables/}%
  {Chapter 5 - Functions of Random Variables/}%
  {Chapter 6 - Two- and Higher-Dimensional Random Variables/}%
  {Chapter 7 - Further Characteristics of Random Variables/}%
  {Chapter 8 - The Poisson and Other Discrete Random Variables/}%
  {Chapter 9 - Some Important Continuous Random Variables/}%
  {Chapter 10 - The Moment-Generating Function/}%
}
\usepackage{setspace}
\setstretch{0.5}

\usepackage{lmodern}
\usepackage{amsmath,amssymb,amsthm,enumitem,mathtools,xpatch}
\usepackage{bm}
\usepackage[most]{tcolorbox}
\usepackage[dvipsnames]{xcolor}
\usepackage{float}
\usepackage{mathrsfs}
\usepackage{wrapfig, lipsum, amsthm, thmtools}
\usepackage{framed}
\usepackage{gensymb}
\usepackage{url}

\newcommand*{\simsym}{\mathord\sim}

\newcommand*{\Perm}[2]{{}^{#1}\!P_{#2}}%
\newcommand*{\Comb}[2]{{}^{#1}C_{#2}}%

\usepackage[framemethod=tikz]{mdframed}

% ---------------------------------------------------------------------------
% Theorem environments
% ---------------------------------------------------------------------------
\theoremstyle{definition}
\newtheorem*{exmp*}{Example}

% Family used by chapters 1-6, 9, 10
\newtheorem*{defn}{Definition}
\surroundwithmdframed[backgroundcolor=white]{defn}

\newtheorem{cor}{Corollary}
\surroundwithmdframed[backgroundcolor=white]{cor}

\newtheorem{prop}{Proposition}
\surroundwithmdframed[backgroundcolor=white]{prop}

\newtheorem*{thm}{Theorem}
\surroundwithmdframed[backgroundcolor=white]{thm}

% Family used by chapters 7 and 8
\newtheorem*{defn*}{Definition}
\surroundwithmdframed[backgroundcolor=white]{defn*}

\newtheorem{cor*}{Corollary}
\surroundwithmdframed[backgroundcolor=white]{cor*}

\newtheorem*{prop*}{Proposition}
\surroundwithmdframed[backgroundcolor=white]{prop*}

\newtheorem*{thm*}{Theorem}
\surroundwithmdframed[backgroundcolor=white]{thm*}

% ---------------------------------------------------------------------------
% Encoding
% ---------------------------------------------------------------------------
\usepackage[utf8]{inputenc}

% ---------------------------------------------------------------------------
% tikz (probability trees) and pgfplots
% ---------------------------------------------------------------------------
\usepackage{tikz}
\usetikzlibrary{trees,calc,angles,positioning,intersections}

\usepackage{pgfplots}
\pgfplotsset{width=10cm,compat=1.9}
\usepgfplotslibrary{fillbetween}

% quotations dirty talk
\usepackage{dirtytalk}

% floor ceiling
\DeclarePairedDelimiter\ceil{\lceil}{\rceil}
\DeclarePairedDelimiter\floor{\lfloor}{\rfloor}
\DeclarePairedDelimiter\abs{\lvert}{\rvert}
% plain \abs{x} autosizes like \left| x \right|; \abs*{x} keeps fixed-size bars
\DeclarePairedDelimiter\absfixed{\lvert}{\rvert}%
\RenewDocumentCommand{\abs}{s}{\IfBooleanTF{#1}{\absfixed}{\absfixed*}}%

% diag table
\usepackage{diagbox}

% math circle
\makeatletter
\newcommand\mathcircled[1]{%
  \mathpalette\@mathcircled{#1}%
}
\newcommand\@mathcircled[2]{%
  \tikz[baseline=(math.base)] \node[draw,circle,inner sep=1pt] (math) {$\m@th#1#2$};%
}
\makeatother

% ---------------------------------------------------------------------------
% Title block
% ---------------------------------------------------------------------------
% Follows the convention used in the deqla and apde manuals: the document
% defines \nbook (and \nbookshort) before inputting this header, and the title
% block is assembled here from those macros -- \title is left empty and the
% book, its author and the solutions author are stacked in \author, with a
% closing disclaimer appended to \maketitle.
%
% Guarded on \nbook so only the merged document picks this up; the ten chapter
% documents declare their own \title/\author after inputting this file and are
% left exactly as they were. Defining \nisofficial suppresses the disclaimer.
\ifx \nbook \undefined
\else
  \ifx \nauthor \undefined
    \def\nauthor{David A. Lee}
  \fi
  \author{\nbook \\ \small Paul L. Meyer \\ \small Solutions by \nauthor}
  \date{}
  \title{}
  \ifx \nisofficial \undefined
    \makeatletter
    \let\@real@maketitle\maketitle
    \renewcommand{\maketitle}{\@real@maketitle\begin{center}\begin{minipage}[c]{0.9\textwidth}\centering\footnotesize All errors, typographical and substantive, and other offenses, are entirely my own.\end{minipage}\end{center}}
    \makeatother
  \fi
\fi
}{\input{../headermeyer-prob-stats}}
%  (the fallback lets a chapter compile whether the working directory is the
%   repository root or the chapter's own folder).
%
%  This file is the *union* of the ten original per-chapter preambles, so every
%  chapter keeps the packages, macros and theorem environments it relied on.
%  Two divergences in the originals are reconciled here:
%    * inputenc: chapter 1 declared utf8, the others latin1. Every source file
%      is in fact UTF-8, and the only non-ASCII byte in any of them is a "§"
%      inside a comment, so utf8 is used throughout and renders identically.
%    * theorem environments: chapters 7 and 8 used the starred names defn*,
%      cor*, prop*, thm*; the other chapters used defn, cor, prop, thm. Both
%      families are defined below, with their original numbering preserved.
% ============================================================================

\PassOptionsToPackage{dvipsnames}{xcolor}

\usepackage{geometry}
\geometry{letterpaper}
\geometry{
 a4paper,
 total={170mm,257mm},
 left=15mm,
 right = 15mm,
 top=15mm,
 bottom = 20mm
 }

\usepackage{graphicx}
% Figures live alongside the chapter that uses them. Listing "./" first plus
% every chapter folder lets a document find its figures whether it is compiled
% from inside its own folder or from the repository root (which is how the
% master builds), so the chapters no longer depend on the working directory.
\graphicspath{%
  {./}%
  {Chapter 1 - Introduction to Probability/}%
  {Chapter 2 - Finite Sample Spaces/}%
  {Chapter 3 - Conditional Probability and Independence/}%
  {Chapter 4 - One-Dimensional Random Variables/}%
  {Chapter 5 - Functions of Random Variables/}%
  {Chapter 6 - Two- and Higher-Dimensional Random Variables/}%
  {Chapter 7 - Further Characteristics of Random Variables/}%
  {Chapter 8 - The Poisson and Other Discrete Random Variables/}%
  {Chapter 9 - Some Important Continuous Random Variables/}%
  {Chapter 10 - The Moment-Generating Function/}%
}
\usepackage{setspace}
\setstretch{0.5}

\usepackage{lmodern}
\usepackage{amsmath,amssymb,amsthm,enumitem,mathtools,xpatch}
\usepackage{bm}
\usepackage[most]{tcolorbox}
\usepackage[dvipsnames]{xcolor}
\usepackage{float}
\usepackage{mathrsfs}
\usepackage{wrapfig, lipsum, amsthm, thmtools}
\usepackage{framed}
\usepackage{gensymb}
\usepackage{url}

\newcommand*{\simsym}{\mathord\sim}

\newcommand*{\Perm}[2]{{}^{#1}\!P_{#2}}%
\newcommand*{\Comb}[2]{{}^{#1}C_{#2}}%

\usepackage[framemethod=tikz]{mdframed}

% ---------------------------------------------------------------------------
% Theorem environments
% ---------------------------------------------------------------------------
\theoremstyle{definition}
\newtheorem*{exmp*}{Example}

% Family used by chapters 1-6, 9, 10
\newtheorem*{defn}{Definition}
\surroundwithmdframed[backgroundcolor=white]{defn}

\newtheorem{cor}{Corollary}
\surroundwithmdframed[backgroundcolor=white]{cor}

\newtheorem{prop}{Proposition}
\surroundwithmdframed[backgroundcolor=white]{prop}

\newtheorem*{thm}{Theorem}
\surroundwithmdframed[backgroundcolor=white]{thm}

% Family used by chapters 7 and 8
\newtheorem*{defn*}{Definition}
\surroundwithmdframed[backgroundcolor=white]{defn*}

\newtheorem{cor*}{Corollary}
\surroundwithmdframed[backgroundcolor=white]{cor*}

\newtheorem*{prop*}{Proposition}
\surroundwithmdframed[backgroundcolor=white]{prop*}

\newtheorem*{thm*}{Theorem}
\surroundwithmdframed[backgroundcolor=white]{thm*}

% ---------------------------------------------------------------------------
% Encoding
% ---------------------------------------------------------------------------
\usepackage[utf8]{inputenc}

% ---------------------------------------------------------------------------
% tikz (probability trees) and pgfplots
% ---------------------------------------------------------------------------
\usepackage{tikz}
\usetikzlibrary{trees,calc,angles,positioning,intersections}

\usepackage{pgfplots}
\pgfplotsset{width=10cm,compat=1.9}
\usepgfplotslibrary{fillbetween}

% quotations dirty talk
\usepackage{dirtytalk}

% floor ceiling
\DeclarePairedDelimiter\ceil{\lceil}{\rceil}
\DeclarePairedDelimiter\floor{\lfloor}{\rfloor}
\DeclarePairedDelimiter\abs{\lvert}{\rvert}
% plain \abs{x} autosizes like \left| x \right|; \abs*{x} keeps fixed-size bars
\DeclarePairedDelimiter\absfixed{\lvert}{\rvert}%
\RenewDocumentCommand{\abs}{s}{\IfBooleanTF{#1}{\absfixed}{\absfixed*}}%

% diag table
\usepackage{diagbox}

% math circle
\makeatletter
\newcommand\mathcircled[1]{%
  \mathpalette\@mathcircled{#1}%
}
\newcommand\@mathcircled[2]{%
  \tikz[baseline=(math.base)] \node[draw,circle,inner sep=1pt] (math) {$\m@th#1#2$};%
}
\makeatother

% ---------------------------------------------------------------------------
% Title block
% ---------------------------------------------------------------------------
% Follows the convention used in the deqla and apde manuals: the document
% defines \nbook (and \nbookshort) before inputting this header, and the title
% block is assembled here from those macros -- \title is left empty and the
% book, its author and the solutions author are stacked in \author, with a
% closing disclaimer appended to \maketitle.
%
% Guarded on \nbook so only the merged document picks this up; the ten chapter
% documents declare their own \title/\author after inputting this file and are
% left exactly as they were. Defining \nisofficial suppresses the disclaimer.
\ifx \nbook \undefined
\else
  \ifx \nauthor \undefined
    \def\nauthor{David A. Lee}
  \fi
  \author{\nbook \\ \small Paul L. Meyer \\ \small Solutions by \nauthor}
  \date{}
  \title{}
  \ifx \nisofficial \undefined
    \makeatletter
    \let\@real@maketitle\maketitle
    \renewcommand{\maketitle}{\@real@maketitle\begin{center}\begin{minipage}[c]{0.9\textwidth}\centering\footnotesize All errors, typographical and substantive, and other offenses, are entirely my own.\end{minipage}\end{center}}
    \makeatother
  \fi
\fi
}{\input{../headermeyer-prob-stats}}
%  (the fallback lets a chapter compile whether the working directory is the
%   repository root or the chapter's own folder).
%
%  This file is the *union* of the ten original per-chapter preambles, so every
%  chapter keeps the packages, macros and theorem environments it relied on.
%  Two divergences in the originals are reconciled here:
%    * inputenc: chapter 1 declared utf8, the others latin1. Every source file
%      is in fact UTF-8, and the only non-ASCII byte in any of them is a "§"
%      inside a comment, so utf8 is used throughout and renders identically.
%    * theorem environments: chapters 7 and 8 used the starred names defn*,
%      cor*, prop*, thm*; the other chapters used defn, cor, prop, thm. Both
%      families are defined below, with their original numbering preserved.
% ============================================================================

\PassOptionsToPackage{dvipsnames}{xcolor}

\usepackage{geometry}
\geometry{letterpaper}
\geometry{
 a4paper,
 total={170mm,257mm},
 left=15mm,
 right = 15mm,
 top=15mm,
 bottom = 20mm
 }

\usepackage{graphicx}
% Figures live alongside the chapter that uses them. Listing "./" first plus
% every chapter folder lets a document find its figures whether it is compiled
% from inside its own folder or from the repository root (which is how the
% master builds), so the chapters no longer depend on the working directory.
\graphicspath{%
  {./}%
  {Chapter 1 - Introduction to Probability/}%
  {Chapter 2 - Finite Sample Spaces/}%
  {Chapter 3 - Conditional Probability and Independence/}%
  {Chapter 4 - One-Dimensional Random Variables/}%
  {Chapter 5 - Functions of Random Variables/}%
  {Chapter 6 - Two- and Higher-Dimensional Random Variables/}%
  {Chapter 7 - Further Characteristics of Random Variables/}%
  {Chapter 8 - The Poisson and Other Discrete Random Variables/}%
  {Chapter 9 - Some Important Continuous Random Variables/}%
  {Chapter 10 - The Moment-Generating Function/}%
}
\usepackage{setspace}
\setstretch{0.5}

\usepackage{lmodern}
\usepackage{amsmath,amssymb,amsthm,enumitem,mathtools,xpatch}
\usepackage{bm}
\usepackage[most]{tcolorbox}
\usepackage[dvipsnames]{xcolor}
\usepackage{float}
\usepackage{mathrsfs}
\usepackage{wrapfig, lipsum, amsthm, thmtools}
\usepackage{framed}
\usepackage{gensymb}
\usepackage{url}

\newcommand*{\simsym}{\mathord\sim}

\newcommand*{\Perm}[2]{{}^{#1}\!P_{#2}}%
\newcommand*{\Comb}[2]{{}^{#1}C_{#2}}%

\usepackage[framemethod=tikz]{mdframed}

% ---------------------------------------------------------------------------
% Theorem environments
% ---------------------------------------------------------------------------
\theoremstyle{definition}
\newtheorem*{exmp*}{Example}

% Family used by chapters 1-6, 9, 10
\newtheorem*{defn}{Definition}
\surroundwithmdframed[backgroundcolor=white]{defn}

\newtheorem{cor}{Corollary}
\surroundwithmdframed[backgroundcolor=white]{cor}

\newtheorem{prop}{Proposition}
\surroundwithmdframed[backgroundcolor=white]{prop}

\newtheorem*{thm}{Theorem}
\surroundwithmdframed[backgroundcolor=white]{thm}

% Family used by chapters 7 and 8
\newtheorem*{defn*}{Definition}
\surroundwithmdframed[backgroundcolor=white]{defn*}

\newtheorem{cor*}{Corollary}
\surroundwithmdframed[backgroundcolor=white]{cor*}

\newtheorem*{prop*}{Proposition}
\surroundwithmdframed[backgroundcolor=white]{prop*}

\newtheorem*{thm*}{Theorem}
\surroundwithmdframed[backgroundcolor=white]{thm*}

% ---------------------------------------------------------------------------
% Encoding
% ---------------------------------------------------------------------------
\usepackage[utf8]{inputenc}

% ---------------------------------------------------------------------------
% tikz (probability trees) and pgfplots
% ---------------------------------------------------------------------------
\usepackage{tikz}
\usetikzlibrary{trees,calc,angles,positioning,intersections}

\usepackage{pgfplots}
\pgfplotsset{width=10cm,compat=1.9}
\usepgfplotslibrary{fillbetween}

% quotations dirty talk
\usepackage{dirtytalk}

% floor ceiling
\DeclarePairedDelimiter\ceil{\lceil}{\rceil}
\DeclarePairedDelimiter\floor{\lfloor}{\rfloor}
\DeclarePairedDelimiter\abs{\lvert}{\rvert}
% plain \abs{x} autosizes like \left| x \right|; \abs*{x} keeps fixed-size bars
\DeclarePairedDelimiter\absfixed{\lvert}{\rvert}%
\RenewDocumentCommand{\abs}{s}{\IfBooleanTF{#1}{\absfixed}{\absfixed*}}%

% diag table
\usepackage{diagbox}

% math circle
\makeatletter
\newcommand\mathcircled[1]{%
  \mathpalette\@mathcircled{#1}%
}
\newcommand\@mathcircled[2]{%
  \tikz[baseline=(math.base)] \node[draw,circle,inner sep=1pt] (math) {$\m@th#1#2$};%
}
\makeatother

% ---------------------------------------------------------------------------
% Title block
% ---------------------------------------------------------------------------
% Follows the convention used in the deqla and apde manuals: the document
% defines \nbook (and \nbookshort) before inputting this header, and the title
% block is assembled here from those macros -- \title is left empty and the
% book, its author and the solutions author are stacked in \author, with a
% closing disclaimer appended to \maketitle.
%
% Guarded on \nbook so only the merged document picks this up; the ten chapter
% documents declare their own \title/\author after inputting this file and are
% left exactly as they were. Defining \nisofficial suppresses the disclaimer.
\ifx \nbook \undefined
\else
  \ifx \nauthor \undefined
    \def\nauthor{David A. Lee}
  \fi
  \author{\nbook \\ \small Paul L. Meyer \\ \small Solutions by \nauthor}
  \date{}
  \title{}
  \ifx \nisofficial \undefined
    \makeatletter
    \let\@real@maketitle\maketitle
    \renewcommand{\maketitle}{\@real@maketitle\begin{center}\begin{minipage}[c]{0.9\textwidth}\centering\footnotesize All errors, typographical and substantive, and other offenses, are entirely my own.\end{minipage}\end{center}}
    \makeatother
  \fi
\fi
